\documentclass{beamer}

\usetheme{Berlin}
\usecolortheme{Orchid}
\useoutertheme{miniframes}
\setbeamertemplate{navigation symbols}{}

\usepackage[utf8]{inputenc}
\usepackage[T1]{fontenc}
\usepackage{amsmath}
\usepackage{amssymb}
\usepackage{booktabs}
\usepackage{graphicx}
\graphicspath{{../images/}}

\usepackage{xcolor}
\usepackage{listings}
\usepackage{hyperref}
\usepackage{tikz}
\usetikzlibrary{arrows.meta,positioning,shapes.geometric,calc}

\lstset{
  language=Java,
  basicstyle=\ttfamily\small,
  keywordstyle=\color{blue},
  commentstyle=\color{gray},
  stringstyle=\color{teal},
  showstringspaces=false,
  columns=fullflexible,
  breaklines=true
}

% Per-slide source line (references shown on the slide itself).
\newcommand{\slidesources}[1]{%
  \vfill
  {\tiny\color{gray}\raggedright\sloppy\parbox{\linewidth}{Sources: #1}\par}%
}

% Unit-wise source strings for this subject (used by slides/notes).
% Path is relative to the lecture `latex/` folder (the compilation CWD).
% OOPS with Java (CSEG1044) — unit-wise sources
%
% These are short, student-facing reference strings intended to be used with:
%   - \slidesources{...}  (in beamer slides)
%   - \notesources{...}   (in lecture notes)
%
% Style inspiration: other subjects in My-Subjects/ use semicolon-separated sources
% and include an "accessed" date for web documentation.

\newcommand{\OOPSUnitISources}{Schildt, \textit{Java: A Beginner's Guide} (10e), McGraw-Hill (Java fundamentals + OOP basics).; Oracle Java Tutorials: Getting Started + Language Basics + Classes and Objects (accessed \today).; Java SE API docs (e.g., \texttt{String}, \texttt{Scanner}) (accessed \today).}

\newcommand{\OOPSUnitIISources}{Schildt, \textit{Java: The Complete Reference} (12e), McGraw-Hill (classes, inheritance, packages, interfaces).; Bloch, \textit{Effective Java} (3e), Addison-Wesley, 2018 (design choices; interfaces vs abstract classes).; Oracle Java Tutorials: Classes and Objects + Interfaces + Packages (accessed \today).; Java SE API docs (e.g., \texttt{Comparable}, \texttt{Runnable}) (accessed \today).}

\newcommand{\OOPSUnitIIISources}{Schildt, \textit{Java: The Complete Reference} (12e) (nested/inner classes, exceptions, threads, I/O).; Oracle Java Tutorials: Nested Classes + Exceptions + Concurrency + Basic I/O (accessed \today).; Java SE API docs (e.g., \texttt{Thread}, \texttt{AutoCloseable}) (accessed \today).}

\newcommand{\OOPSUnitIVSources}{Schildt, \textit{Java: The Complete Reference} (12e) (generics, lambdas, Swing, JDBC).; Bloch, \textit{Effective Java} (3e) (generics + lambdas).; Oracle Java Tutorials: Generics + Lambda Expressions + Swing + JDBC Basics (accessed \today).; Java SE API docs (e.g., \texttt{java.util.function}, \texttt{javax.swing}, \texttt{java.sql}) (accessed \today).}

\newcommand{\OOPSUnitVSources}{Schildt, \textit{Java: The Complete Reference} (12e) (Collections Framework + wrapper classes).; Oracle Java docs: Collections Framework overview (accessed \today).; Java SE API docs (\texttt{java.util} collections; wrapper classes in \texttt{java.lang}) (accessed \today).}

\newcommand{\OOPSUnitVISources}{Git SCM: \textit{Pro Git} (2e) (version control workflow), accessed \today.; GitHub Docs (repositories, issues, pull requests), accessed \today.; Oracle Java Tutorials: Swing + JDBC (accessed \today).; JUnit 5 User Guide (accessed \today).; Apache Maven Project: Getting Started (accessed \today).; Gradle User Manual: Getting Started (accessed \today).; UML class diagrams: UML 2.x overview/spec (accessed \today).}

\newcommand{\OOPSMiscWhyInterfacesSources}{Schildt, \textit{Java: The Complete Reference} (12e) (interfaces).; Bloch, \textit{Effective Java} (3e) (prefer interfaces where appropriate).; Oracle Java Tutorials: Interfaces (accessed \today).; Java SE API docs (e.g., \texttt{Comparable}, \texttt{Runnable}) (accessed \today).}



\title[OOPS with Java]{Object Oriented Programming (CSEG1044)}
\subtitle{Unit 04 -- Lecture 04: Swing Basics (Components, Containers, Layouts)}
\begin{document}

\begin{frame}[fragile]
  \titlepage
  \vspace{0.5em}
  \slidesources{\OOPSUnitIVSources}
\end{frame}

\begin{frame}{Quick Links}
  \centering
  \hyperlink{sec:overview}{\beamerbutton{Overview}}\hspace{0.6em}
  \hyperlink{sec:concepts}{\beamerbutton{Key Topics}}\hspace{0.6em}
  \hyperlink{sec:demo}{\beamerbutton{Demo}}\hspace{0.6em}
  \hyperlink{sec:summary}{\beamerbutton{Summary}}
  \slidesources{\OOPSUnitIVSources}
\end{frame}

\begin{frame}{Agenda}
  \tableofcontents
  \slidesources{\OOPSUnitIVSources}
\end{frame}

\section{Overview}
\label{sec:overview}

\begin{frame}{Lecture Snapshot}
\begin{itemize}[<+->]
  \item \textbf{Unit focus:} Generics, Lambdas, Swing, and JDBC
  \item \textbf{Course thread:} Use Swing Basics (Components, Containers, Layouts) to move Campus Resource Manager toward a real desktop application with UI, callbacks, and persistence.
\end{itemize}
  \slidesources{\OOPSUnitIVSources}
\end{frame}

\begin{frame}{Recall Before You Start}
\begin{itemize}[<+->]
  \item \textbf{Classes and packages}: Swing applications still depend on clear classes and package boundaries. Main reference: \texttt{Unit 02 / Lecture 03 slides/notes}.
  \item \textbf{Nested helpers}: GUI code often uses callback and helper patterns introduced earlier. Main reference: \texttt{Unit 03 / Lecture 01 slides/notes}.
\end{itemize}
  \slidesources{\OOPSUnitIVSources}
\end{frame}


\begin{frame}{Learning Outcomes}
  \begin{itemize}[<+->]
    \item Explain Swing component model: \texttt{JFrame}, \texttt{JPanel}, components vs containers.
    \item Use layout managers (Flow/Border/Grid) to structure a UI.
    \item Build and run a small multi-panel UI skeleton.
  \end{itemize}
  \slidesources{\OOPSUnitIVSources}
\end{frame}

\begin{frame}{Key Terms to Know}
\begin{itemize}[<+->]
  \item \textbf{Component}: UI element like button, label, text field.
  \item \textbf{Container}: holds components (e.g., \texttt{JPanel}, \texttt{JFrame} content pane).
  \item \textbf{Layout manager}: decides how components are arranged (sizes/positions).
  \item \textbf{Event Dispatch Thread (EDT)}: Swing UI thread (advanced but important for correctness).
  \item \textbf{MVC idea}: separate UI from business logic (we will use this later with JDBC).
\end{itemize}
  \slidesources{\OOPSUnitIVSources}
\end{frame}

\begin{frame}{Study Flow for This Lecture}
\begin{enumerate}[<+->]
  \item Refresh the prerequisite bridge before reading new ideas in isolation.
  \item Study the main build in this order: \textit{Swing overview: \texttt{JFrame}, \texttt{JPanel}, components vs containers} $\rightarrow$ \textit{Layout managers (Flow/Border/Grid)} $\rightarrow$ \textit{UI composition: multi-panel skeleton}.
  \item Trace the worked example and connect each line back to one concept frame.
  \item Attempt the mini exercises before moving to the recap and exit check.
\end{enumerate}
  \slidesources{\OOPSUnitIVSources}
\end{frame}


\section{Key Topics}
\label{sec:concepts}

\begin{frame}{Unit 04: Generics, Lambdas, GUI Swing \& Database Connectivity}
  \begin{itemize}[<+->]
    \item Lecture focus: Swing Basics (Components, Containers, Layouts)
  \end{itemize}
  \slidesources{\OOPSUnitIVSources}
\end{frame}

\begin{frame}{Today: Roadmap}
  \begin{itemize}[<+->]
    \item Swing overview; `JFrame`, `JPanel`, components vs containers
    \item Layout managers (Flow/Border/Grid) and UI composition
    \item Building a small multi-panel UI skeleton
  \end{itemize}
  \slidesources{\OOPSUnitIVSources}
\end{frame}

\begin{frame}{Components vs containers}
  \begin{itemize}[<+->]
    \item Components: \texttt{JLabel}, \texttt{JButton}, \texttt{JTextField}, \texttt{JTable}, ...
    \item Containers: \texttt{JPanel}, \texttt{JFrame} content pane
    \item Compose UI by nesting panels and adding components inside them.
  \end{itemize}
  \slidesources{\OOPSUnitIVSources}
\end{frame}

\begin{frame}{Layout managers (quick)}
  \begin{itemize}[<+->]
    \item \texttt{FlowLayout}: left-to-right (good for button rows)
    \item \texttt{BorderLayout}: NORTH/SOUTH/EAST/WEST/CENTER (good for screens)
    \item \texttt{GridLayout}: rows/cols (good for simple forms)
  \end{itemize}
  \slidesources{\OOPSUnitIVSources}
\end{frame}

\begin{frame}{A common UI composition pattern}
  \begin{itemize}[<+->]
    \item Root panel: \texttt{BorderLayout}
    \item NORTH: title / header
    \item CENTER: main panel (form/table)
    \item SOUTH: action buttons
  \end{itemize}
  \slidesources{\OOPSUnitIVSources}
\end{frame}

\begin{frame}{Checkpoint and Reflection}
\begin{block}{Reflection prompt}
Where would \textit{Swing Basics (Components, Containers, Layouts)} matter most in Campus Resource Manager, and which design choice would you justify using this lecture's ideas?
\end{block}
\begin{block}{Quick self-check}
Which part needs one more example before you move on: \textit{Swing overview: \texttt{JFrame}, \texttt{JPanel}, components vs containers}, \textit{Layout managers (Flow/Border/Grid)}, the worked example, or the recap?
\end{block}
  \slidesources{\OOPSUnitIVSources}
\end{frame}

\section{Demo / Example}
\label{sec:demo}

\begin{frame}[fragile]{Example: multi-panel skeleton (1/3)}
\begin{lstlisting}
import javax.swing.*;
import java.awt.*;

public class SimpleFormUI {
  public static void main(String[] args) {
    SwingUtilities.invokeLater(() -> {
      JFrame f = new JFrame("CRM Skeleton");
      f.setDefaultCloseOperation(JFrame.EXIT_ON_CLOSE);

      JPanel root = new JPanel(new BorderLayout(10, 10));
\end{lstlisting}
  \slidesources{\OOPSUnitIVSources}
\end{frame}

\begin{frame}[fragile]{Example: multi-panel skeleton (2/3)}
\begin{lstlisting}
// continued...
      root.add(new JLabel("Add Resource", SwingConstants.CENTER),
               BorderLayout.NORTH);

      JPanel form = new JPanel(new GridLayout(2, 2, 8, 8));
      form.add(new JLabel("Name:"));
      form.add(new JTextField());
      form.add(new JLabel("Type:"));
      form.add(new JTextField());
      root.add(form, BorderLayout.CENTER);
\end{lstlisting}
  \slidesources{\OOPSUnitIVSources}
\end{frame}

\begin{frame}[fragile,shrink=8]{Example: multi-panel skeleton (3/3)}
\begin{lstlisting}
// continued...
      JPanel actions = new JPanel(new FlowLayout(FlowLayout.RIGHT));
      actions.add(new JButton("Save"));
      actions.add(new JButton("Exit"));
      root.add(actions, BorderLayout.SOUTH);

      f.setContentPane(root);
      f.pack();
      f.setLocationRelativeTo(null);
      f.setVisible(true);
    });
  }
}
\end{lstlisting}
  \slidesources{\OOPSUnitIVSources}
\end{frame}

\begin{frame}{Mini Exercises (2--3 mins)}
  \begin{enumerate}[<+->]
    \item Which is a container: \texttt{JPanel} or \texttt{JButton}?
    \item Which layout is best for header + body + footer screens?
    \item Why prefer \texttt{pack()} over fixed \texttt{setSize()} (often)?
  \end{enumerate}
  \slidesources{\OOPSUnitIVSources}
\end{frame}

\begin{frame}{Watch-outs}
\begin{itemize}[<+->]
  \item Putting every control directly on one JFrame without structure.
  \item Hardcoding coordinates instead of using layout managers intentionally.
  \item Mixing UI code with business logic too early.
\end{itemize}
  \slidesources{\OOPSUnitIVSources}
\end{frame}

\section{Summary}
\label{sec:summary}

\begin{frame}{Key Takeaways}
  \begin{itemize}[<+->]
    \item Swing UI = components inside containers.
    \item Layout managers create robust UIs without hard-coded positions.
    \item Multi-panel composition scales better than one giant panel.
  \end{itemize}
  \slidesources{\OOPSUnitIVSources}
\end{frame}

\begin{frame}{Checkpoint Questions}
  \begin{enumerate}[<+->]
    \item What is the difference between a component and a container?
    \item Name 3 layout managers and one use case for each.
    \item Describe one UI screen using the NORTH/CENTER/SOUTH pattern.
  \end{enumerate}
  \slidesources{\OOPSUnitIVSources}
\end{frame}

\begin{frame}{Next Study Steps}
\begin{itemize}[<+->]
  \item Revisit the recall references if any older idea still feels weak.
  \item Rework the lecture example without looking at the notes first.
  \item Attempt the self-study practice and then answer the exit questions in full sentences.
  \item Apply the idea once inside Campus Resource Manager to make the concept stick.
\end{itemize}
  \slidesources{\OOPSUnitIVSources}
\end{frame}

\begin{frame}{Next Lecture Preview}
  \textbf{Next:} Swing event handling + JDBC (persistence)
  \vspace{0.6em}
  \begin{itemize}[<+->]
    \item We will connect buttons to actions and store data in a database.
  \end{itemize}
  \slidesources{\OOPSUnitIVSources}
\end{frame}

\end{document}
